\section{Experimental Methodology}

\subsection{Research Questions}
The experimental protocol is organized around four research questions. \textbf{RQ1} asks whether Spec2Testbench can execute real and reproducible verification runs on the 28-circuit benchmark used in the paper. \textbf{RQ2} asks whether the framework distinguishes simulability from specification compliance rather than treating successful ngspice execution as sufficient evidence of correctness. \textbf{RQ3} asks whether controlled violations are detected without false accepts or false rejects under the executed mutation campaign. \textbf{RQ4} asks what measurable benefit is provided by the LLM-assisted layer relative to the deterministic baseline. No \textbf{RQ5} on expert agreement is introduced here because no completed human-validation study was located in the canonical evidence audited for this manuscript.

\subsection{Benchmark and Evaluation Scope}
The nominal benchmark evaluated in the paper contains 28 local netlists listed in \texttt{benchmark/analogcoder\_pro/} and referenced by the benchmark manifests and specification files. The repository metadata repeatedly labels this corpus as \texttt{AnalogCoder-Pro}, but the manuscript should describe the local evaluation set as \emph{benchmark-aligned} rather than as a formally provenance-verified mirror of an official upstream release, because the present audit did not locate a cryptographic or publication-grade provenance chain linking every local netlist to an official external benchmark package. The benchmark layout documentation in \texttt{benchmark/README.md} identifies \texttt{benchmark/analogcoder\_pro} as the current pedagogical benchmark tier, while the example specification files label the technology as \texttt{AnalogCoder-Pro/PySpice generic Level-1 models}. This wording is the safe basis for the manuscript.

At the level of executed nominal evidence, the 28 circuits span 14 circuit families: amplifier, current mirror, comparator, low-pass filter, high-pass filter, band-pass filter, notch filter, opamp, mixer, oscillator, opamp integrator, opamp differentiator, composite arithmetic blocks, and Schmitt trigger. The local netlists are compact rather than industrial-scale: a direct audit of the 28 \texttt{.cir} files in \texttt{benchmark/analogcoder\_pro/} gives a range of 5 to 24 non-empty lines per netlist, with a median of 13 non-empty lines, and a range of 3 to 18 statement-like lines, with a median of 9. These counts are included only as a descriptive complexity proxy for the local benchmark files and should not be overstated as a universal circuit-complexity metric.

The benchmark manifests also document the analyses and expected metrics available to the framework. Across the 28 cases, the configured analyses include DC, AC, transient, and Fourier/spectral paths, depending on circuit family. The configured metric space includes steady-state quantities such as operating point and quiescent current, AC quantities such as \texttt{dc\_gain\_db}, cutoff frequency, and phase margin, transient quantities such as propagation delay, slew-related behavior, settling time, oscillator startup amplitude, and oscillator frequency, and spectral quantities such as THD and fundamental frequency. In the canonical paper campaign, the actually extracted nominal evidence is narrower than the full manifest capability: \texttt{results/paper\_metric\_results.csv} contains executed nominal measurements over operating point, quiescent current, \texttt{dc\_gain\_db}, phase margin, cutoff frequency, propagation delay, oscillator frequency, startup amplitude, slew rate, settling time, THD, and fundamental frequency.

\subsection{Execution Environment}
The canonical environment record is \texttt{results/final\_environment.json}. It reports Windows 10 as the operating system, Python 3.12.2, ngspice \texttt{ngspice-41 : Circuit level simulation program}, architecture \texttt{AMD64}, and Git commit \texttt{2373572107fae2e602abf0d0f918be081af09c89}. The same record reports \texttt{mock\_allowed = false} for the audited campaign environment and \texttt{pyspice\_disabled\_campaign = true}. The canonical native measurement backends are \texttt{NGSPICE\_MEASURE} and \texttt{NGSPICE\_WRDATA}; both are supported by direct executed evidence, but their roles differ across campaigns. \texttt{NGSPICE\_MEASURE} is the dominant backend in the nominal and frozen-pilot evidence, whereas \texttt{NGSPICE\_WRDATA} is validated by the two-case extension with independent comparisons.

The paper campaign configuration itself is recorded in \texttt{configs/paper\_experiment.yaml}. It specifies a reproducible nominal campaign with \texttt{paper\_eligible\_only: true}, \texttt{allow\_mock: false}, raw-output and provenance preservation enabled, and \texttt{llm.enabled: false}. This point is methodologically important: the canonical nominal paper campaign is a deterministic run without an active LLM generation path. As a consequence, no canonical manuscript evidence supports concrete values for LLM temperature, top-p, top-k, candidate count, or random seeds for the executed paper campaign. If a summary table includes these fields, they should be filled with \texttt{TBD -- TO BE FILLED FROM CANONICAL EVIDENCE} or marked as not applicable for the executed nominal run. The same caution applies to framework version strings when they are not explicitly recorded in the canonical environment artifact.

\subsection{Experimental Campaigns}
The manuscript should distinguish six campaign families and report them separately. The first is the nominal ACP-28 campaign, for which the canonical evidence is \texttt{results/paper\_campaign\_summary.json} and \texttt{results/paper\_campaign\_summary.csv}. This campaign executed 28 benchmark-aligned circuits in \texttt{REAL} mode, with \texttt{execution\_status = SUCCESS} for all cases, \texttt{paper\_eligible = True} for all cases, 27 \texttt{SIMULABLE\_COMPLIANT} outcomes, and one \texttt{SIMULABLE\_NONCOMPLIANT} outcome. The second campaign family is the controlled-violation study. The canonical current summary is \texttt{results/final\_controlled\_summary.json}, which reports 30 generated variants but only 2 effective controlled violations, with one detected effective violation and one false pass. This campaign therefore supports a limited statement about executed controlled violations, but it does not support a blanket claim of zero false accepts or zero false rejects.

The third campaign family is backend validation. This is supported by the frozen pilot V3 and WRDATA extension evidence rather than by the software test suite alone. The canonical frozen pilot summary \texttt{results/frozen\_pilot\_metrics\_v3.json} records 16 cases with \texttt{TRUE\_ACCEPT = 8}, \texttt{TRUE\_DETECTION = 8}, and no false accepts, false rejects, or unevaluated cases in that frozen dataset. The WRDATA validation evidence is recorded in \texttt{results/wrdata\_independent\_comparisons.csv}, where both independently recomputed vector comparisons agree within tolerance. These backend experiments evaluate evidence extraction behavior and should not be numerically merged with the nominal ACP-28 campaign or with the later expanded controlled-violation campaign.

The fourth, fifth, and sixth campaign families are currently incomplete or absent as scientific evidence. The LLM-versus-deterministic ablation remains partial: \texttt{results/final\_ablation\_summary.json} records \texttt{status = PARTIAL}, \texttt{A2 = NOT\_EXECUTED}, \texttt{A4 = NOT\_EXECUTED}, and \texttt{llm\_included = false}. Variation-aware checks and robustness remain unexecuted at the level required for quantitative claims, since \texttt{results/final\_robustness\_metrics.json} records \texttt{status = NOT\_EXECUTED}. Expert validation likewise remains unavailable in the canonical evidence used here. These three campaign families must therefore be presented as future work, limitations, or \texttt{TBD} entries rather than as completed experimental blocks.

\subsection{Experimental Metrics}
The reported scientific metrics must be tied to the appropriate campaign and must not be confused with software test counts. For the nominal ACP-28 campaign, the relevant metrics are simulation success rate, real simulation rate, paper eligibility rate, nominal compliance rate, simulable-but-noncompliant count, and metric coverage. In the canonical nominal campaign, these quantities are directly available from \texttt{results/paper\_campaign\_summary.json}: simulation success rate is 1.0, real simulation rate is 1.0, paper eligibility rate is 1.0, nominal compliance rate is 0.9642857142857143, simulable-but-noncompliant rate is 0.03571428571428571, and metric extraction success rate is 1.0. Because all 28 nominal cases have nonzero extracted metric counts, the nominal campaign contains no circuits without extracted metrics in the canonical summary.

For the backend-validation and controlled-violation campaigns, the manuscript should use confusion-style metrics that match the actual dataset under discussion. The frozen pilot V3 supports \texttt{TRUE\_ACCEPT}, \texttt{TRUE\_DETECTION}, \texttt{FALSE\_ACCEPT}, \texttt{FALSE\_REJECT}, and \texttt{UNEVALUATED} counts exactly as recorded in \texttt{results/frozen\_pilot\_metrics\_v3.json}. The larger controlled-violation campaign supports effective-violation recall, false-pass rate, decision coverage, and unevaluated rate only in its own restricted scope, as recorded in \texttt{results/final\_controlled\_summary.json}. These counts must not be merged into a single unified table unless the table explicitly separates datasets.

No canonical executed evidence currently supports a quantitative estimate of LLM cost, LLM iteration count, or LLM sampling hyperparameters for the paper campaign, because the campaign configuration disables the LLM and the ablation artifact remains partial. If the experimental-methodology section includes a compact campaign table, the LLM-specific fields should therefore be marked \texttt{TBD -- TO BE FILLED FROM CANONICAL EVIDENCE} or \texttt{not executed}. Likewise, provider unavailability must not be conflated with circuit failure: LLM configuration gaps belong to the method boundary and ablation limitations, whereas circuit execution outcomes belong to the ngspice-based scientific campaigns.
